\documentclass[DIV=15, fleqn, oneside, 10pt]{scrartcl}

% --------------------------------------------------
% Page appearance (subtle warm tone)
% --------------------------------------------------
\usepackage{xcolor}
\pagecolor{yellow!3}        
\color{black}    

% --------------------------------------------------
% Layout and spacing
% --------------------------------------------------
\KOMAoptions{parskip=half}  % space between paragraphs (no indentation feel)
\setlength{\parindent}{0pt} % remove paragraph indentation
\setlength{\mathindent}{0pt}% left-align equations (fleqn tweak)

% --------------------------------------------------
% Math packages
% --------------------------------------------------
\usepackage{amsmath, amssymb}
\AtBeginDocument{\setlength{\mathindent}{0pt}}

% --------------------------------------------------
% Theorem styling (clean, minimal)
% --------------------------------------------------
\usepackage{amsthm, tikz}

\newtheoremstyle{clean}
  {6pt}        % Space above
  {6pt}        % Space below
  {}% Body font
  {}           % Indent amount
  {\bfseries}  % Head font (bold)
  {.}          % Punctuation after heading
  {0.5em}      % Space after heading
  {}

\theoremstyle{clean}
\newtheorem{question}{Q} % auto-numbered questions

% --------------------------------------------------
% Line
% --------------------------------------------------
\newcommand{\sep}{%
  \par\medskip
  \hrule
  \bigskip
}

% --------------------------------------------------
% Title (no date)
% --------------------------------------------------
\title{Interest Rates \& Bonds}
\date{}                    % removes date

% --------------------------------------------------
\begin{document}
\maketitle

\newpage

\section*{What Is An Interest Rate}
\sep
\textbf{Interest is the cost of borrowing or the reward of lending.} Interest rates are used for compounding
when calculating the future value and discounting when used to calculate the present value. The
compounding frequency k, represents the amount of times the rate is applied to compound.

e.g. with interest rate r, compounding frequency k, and time period of t years,
the future and the present value are related
as
\[
V_{F} = V_{P} (1 + \frac{r}{k})^{kt}
\]

If the compounding frequency tends towards, infinity, the case is of continuous compounding:
\[
V_{F} = V_{P} \times e^{rt}
\]

\section*{Bonds}
\sep
\subsection{Zero Coupon Bond}
A zero coupon bond is a contract between the issuer (borrower) and the buyer (lender). Following
are the terms of contract:

 - \textbf{B}: Price of the bond (\$)

 - \textbf{P}: Face value of the bond (\$)

 - \textbf{T}: Time to maturity (yrs)

 - The buyer pays a sum of \textbf{B} to the issuer. The issuer returns a sum of \textbf{P} to
    the buyer after a period of \textbf{T} years from the date of contract.

The interest rate $r$, earned by the buyer is called the \textbf{zero rate} corresponding
to maturity T:
\[
P = B \times e^{rT}
\]

\textbf{Zero rate curve} is the plot of zero rates Vs. maturity for all zero rates available in the
market. Naturally, this plot is not continuous due to the lack of availability of bonds that mature
at any given time T. To make it continous, we interpolate the data. One of the simplest methods is
linear interpolation.

Let $r2$ and $r1$ be the zero rates for maturities $T_2$ and $T_1$ respectively, such that 
$T_1 < T_2$. The zero rate r, for maturity T, between $T_1$ and $T_2$ in this case is given by:
\[
r = r1 + (r2 - r1) \times \frac{T - T_1}{T_2 - T_1}
\]

Drawbacks of linear interpolationa are that the zero curve is not smooth., 
the zero curve has \textit{kinks} at the vertices,
it fails to capture the oscillatory nature of actual zero rates.

Another alternative to linear interpolation is \textbf{cubic spline}. But if the cubic spline
alternates rapidly, it might lead to huge errors in estimation. Also, this method is \textbf{
    globally sensitive
} i.e. changing one data point affects the entire curve. A cubic spline can also sometimes esitmate
negative forward rates which may not ever happen.

The gold standard for plotting zero rates is \textbf{monotone convex interpolation}. It seeks out
a smooth functon $f$ such that the average value between data points $T_i$ and $T_{i + 1}$ is
the discrete forward rate, $F_i$ obtained from linear interpolation:
\[
F_i = \frac{1}{T_i - T_{i-1}} \int_{T_{i-1}}^{T_i} f(t) dt \text{ for all data points i}
\]

\newpage
\subsection{Coupon Paying Bond}
A zero coupon bond is a contract between the issuer (borrower) and the buyer (lender). Following
are the terms of contract:

 - \textbf{B}: Price of the bond (\$)

 - \textbf{P}: Face value of the bond (\$)

 - \textbf{T}: Time to maturity (yrs)

 - \textbf{Coupons}: Periodic coupons paid at fixed dates as per contract: $C_1$ at time $T_1$,
  $C_2$ at time $T_2$ and so on...

The bond price $B$ in this case is calculated as:

\[
B = \sum_{i = 1}^{n}{C_{i} e^{-r_i T_i}}
\]
The summation includes present value of each coupon discounted with the relevant zero rate and the
last term is the present value of the face value which is returned in the end.

\section*{Yield}
\sep
Yield is that internal rate of return, $y$, which makes the PV of all cashflows equal to the bond
price when used for discounting:
\[
B = \sum_{i = 1}^{n}{C_{i} e^{-y T_i}}
\]

The yield of a zero coupon bond is equal to the zero rate.

As seen in the equation, yield and bond price are in negative correlation.

\textbf{Yield of a zero coupon bond is higher than that of a coupon paying bond with the same
maturity and face value.}

Sometimes, each coupon is a fixed percentage of the par value, which is called the 
\textbf{coupon rate}.

Yields are bond specific, while zero rates are market specific. Two bonds with the same maturity
but different coupon rate will have a different yield. Hence, we must take reference of zero rates
instead of yields while counting returns over a time period.


A bond is said to be \textbf{trading at a premium} if $B > P$ ($y < coupon\:rate$)

A bond is said to be \textbf{trading at a discount} if $B < P$ ($y > coupon\:rate$)

A bond is said to be \textbf{trading at par} if $B = P$ ($y = coupon\:rate$)

\section*{Duration}
\sep
\textbf{Duration is the weighted average time, until the present value of all cashflows from a bond
are recieved.} The weight for each time of cashflow is the present value of the cashflow divided
by the bond price. Naturally, the weights all sum back to 1:
\[
D = \sum_{i = 1}^{n}{T_i \left( \frac{C_i e^{-yT_i}}{B} \right)}
\]

\newpage
\section*{Modified Duration}
\sep

For small changes in bond price, it is true that
\[
\Delta B = \frac{dB}{dy} \Delta y
\]
\[
\frac{dB}{dy} = -\sum_{i = 1}^{n}{T_i C_{i} e^{-r_i T_i}} = -BD
\]
Which gives
\begin{equation}
\Delta B = -BD \Delta y
\end{equation}
If we assume a compounding frequency of k instead of continous compounding and repeat the same
above calculation,
\[
\Delta B = -B \left( \frac{D}{1 + \frac{y}{k}} \right) \Delta y
\]
Above value in the bracket is the modified duration, denoted by $D^*$,
\[D^* = 
\frac{D}{1 + \frac{y}{k}}
\]
Modified duration is the negative rate of change of percentage change in bond price per unit change
in yield.

It measures the sensitivity of the bond price wrt change in bond yield.

\section*{Convexity}
\sep
Convexity is the second derivative of the price wrt yield:
\[C = \frac{d^2 B}{dy^2}\]
From Taylor expansion,
\[B(y + \Delta y) = B(y) + \frac{dB}{dy} \Delta y + \frac{1}{2} \frac{d^2 B}{dy^2} (\Delta y)^2\]
\[\implies B(y + \Delta y) - B(y) = -B(y)D \Delta y + \frac{1}{2}C (\Delta y)^2\]
dividing by $B(y)$, we get

\begin{equation}
  \frac{\Delta B}{B} = -D \Delta y + \frac{C}{2B} (\Delta y)^2
\end{equation}
\end{document}
